\documentclass[12pt]{report}
\usepackage{subfigure}
\usepackage[latin1]{inputenc}
\usepackage{amsmath}
\usepackage{amsfonts}
\usepackage{amssymb}
\usepackage{graphicx}
\usepackage{todonotes}

\author{P. Bachant, M. Wosnik, B. Gunawan, V. Neary}
\title{Test Plan: UNH RM2 Tow Tank Experiment}

\begin{document}

\begin{titlepage}
    \centering
    \vfill
    {\bfseries\Large
        RM2 Tow Tank Experimental Test Plan\\
    }  
            \vskip2cm
            P. Bachant$^1$\\
            M. Wosnik$^1$\\
            B. Gunawan$^2$\\
            V. Neary$^2$\\  
    \vfill
    November 7, 2014
    \vfill
    \includegraphics[width=0.33\textwidth]{figures/unhlogo} \\
    \vspace{0.1in}
    \includegraphics[width=0.42\textwidth]{figures/snllogo} \\
    \vspace{0.1in}
    \includegraphics[width=0.35\textwidth]{figures/doelogo}
    \vfill
    \vfill
\end{titlepage}

\tableofcontents

\chapter{Introduction}

The turbine blades of cross-flow turbines experience a large range of angles of
attack as they rotate around the turbine axis of rotation (``cross-flow
turbine'' means the axis of rotation is perpendicular to flow direction---can be
vertical or horizontal).  This range of angles of attack becomes larger as the
tip speed ratio $\lambda=\omega R/U_\infty$  decreases.  The range of angles of
attack is typically sufficiently large that the blades operate under dynamic
stall, which is a Reynolds number dependent phenomenon. Furthermore, the higher
the solidity of a cross-flow turbine, the lower a tip-speed ratio it will
operate at. Marine Hydrokinetic (MHK) cross-flow turbines typically have higher
solidity than cross-flow wind turbines (VAWT), and hence operate at lower
tip-speed ratios. Since MHK turbines operate in a higher density fluid, unsteady
dynamic effects related to the blades? pitching motion and flow curvature also
become more important.

The performance of cross-flow MHK turbines thus depends on both Reynolds number
and solidity (note that these issues are related, since an average blade chord
Reynolds number, $Re_{c,\mathrm{avg}}?\lambda U_\infty c/ \nu$, can be expressed
in terms of tip speed ratio, which is a function of solidity). If numerical
models are validated with physical model data that was obtained at too low
Reynolds numbers, it cannot be determined whether problems with model
predictions are caused by Reynolds number effects, issues related to higher
solidity, or both.  It is uncertain whether numerical models validated with
physical model data obtained at low Reynolds number should be considered
validated at all, since the scale at which the model will be applied is orders
of magnitude larger. One way to overcome this uncertainty is to show that the
scaled physical model test has become Reynolds number independent, and therefore
validation efforts should be relevant at full-scale.

For example, the effects of Reynolds number on average power output was shown to
be significant on the 2 m Sandia Research Darrieus turbine in wind tunnel
testing (Blackwell et al., 1976): Maximum power coefficient,
$C_{P,\mathrm{max}}$, increased with Reynolds number, $Re_c$, along with a shift
of the location of $C_{P,\mathrm{max}}$ toward lower tip speed ratios due to
delayed blade stall. The effects of Reynolds number were quite dramatic over a
relatively small range of $Re_c \approx 1.1 \times 10^5$--$2.9 \times 10^5$.

To date, attempts to validate the CACTUS model have relied on measurements from
SAFL and UNH.  For the SAFL experiments (RM2), the chord Reynolds number,
$Re_c?\sim 10^4$, was below the threshold needed to properly simulate lift and
stall characteristics.  For the UNH experiment (UNH-RVAT, Bachant and Wosnik,
2013), the chord Reynolds number was sufficiently high at $Re_c \approx 2.7
\time 0^5$, but the chord/radius ratio and solidity (13.4\%) created instability
in the free-wake evolution, which caused significant overestimation of power
coefficient.


\section{Objectives}

\begin{enumerate}

	\item Acquire a Reynolds number independent dataset for the DOE RM2 cross-flow
	turbine.

		\begin{enumerate}
			\item Design and build an appropriate turbine model.
		
			\item Acquire full device performance curves (power and drag coefficients) at
			multiple $Re$, to attempt to find convergence.
		  
			\item Once convergence is found, acquire another performance curve at that
			Reynolds number $Re_0$.
		
			\item If time permits, acquire a detailed wake flow map in the near-wake at
			$Re_0$ using acoustic Doppler velocimetry.
		\end{enumerate}
	
	\item Measure the effects of strut drag on RM2 performance.
	
	\begin{enumerate}
		\item Measure the parasitic torque from support struts by rotating the turbine
		in still water. \label{obj-parasitic_torque}
		
		\item Redo task \ref{obj-parasitic_torque} with cylindrical tubes slid over
		the struts to significantly increase drag.
		
		\item Acquire a performance curve at $Re_0$ with the high-drag struts. 
	\end{enumerate}
	
\end{enumerate}


\chapter{Experimental setup and methods}

\section{Turbine model}

The turbine is to be a 1:6 scale model of the DOE RM2 rotor.

\subsection{Geometric requirements}

Turbine geometry is to be scaled from the RM2 design report \cite{Barone2011}.

\begin{table}[ht]
\centering
\begin{tabular}{l|l|l}
   & Full-scale & Model (1:6) \\
\hline 
Diameter (m)   & 6.45 & 1.075 \\ 
Height (m)     & 4.84 & 0.807 \\ 
Blade root chord (m) & 0.400 & 0.067 \\ 
Blade tip chord (m)  & 0.240 & 0.040 \\ 
Blade profile & NACA 0021 & NACA 0021 \\ 
Blade mount & 1/2 chord & 1/2 chord \\ 
Blade pitch (deg.) & 0.0 & 0.0 \\ 
Strut profile & NACA 0021 & NACA 0021 \\ 
Strut chord (m) & 0.360 & 0.060 \\ 
Shaft diameter (m) &  &  \\ 
\end{tabular}
\caption{RM2 turbine geometric parameters.}
\end{table}

The shaft diameter has still not been chosen.


\section{Facility and instrumentation}

Experiments will be performed in the UNH tow/wave tank. The turbine will be
mounted in a frame built from NACA 0020 struts, attached to the tow carriage by
four linear bearings, which transfer all streamwise force to a pair of S-beam
load cells. The turbine shaft RPM will be controlled by a servo motor system,
which allows prescription of the turbine tip speed ratio. The load torque will
be measured by an inline rotary torque transducer and a load cell mounted at a
fixed distance from the servo motor, providing a redundant measurement. Turbine
shaft angle will be measured using the servo drive's emulated encoder output,
set to $10^5$ counts per turbine shaft revolution. Carriage speed, and therefore
inflow velocity will be measured using a linear encoder with 10 $\mu$m
resolution.


\chapter{Research deliverables}

A full performance curve for the turbine will consist of 20--30 tows at varying
tip speed ratio. Each of these tows will produce raw data for the turbine
torque, drag force on the submerged equipment, turbine shaft angle, carriage
speed, and Vectrino velocity measurements. Raw data files will be saved from our
data acquisition device at a sample sate of 2 kHz, from the motion controller at
1 kHz, and from the Vectrino at 200 Hz. All three devices will be triggered by
the motion controller at the beginning of each run so that they are relatively
synchronized.

Tare torque and drag runs will also be performed to measure the shaft bearing
friction torque and turbine mounting frame drag, respectively. These data will
be similar to the turbine performance data, omitting torque measurements for the
tare drag runs and vice versa.

Raw data will be in HDF5 format, processed data in comma separated value (CSV),
and processing code in Python. Raw data files will be uploaded to figshare,
which will then give the dataset a Digital Object Identifier (DOI) and therefore
a persistent URL. Processed data and processing code will be hosted on GitHub to
help facilitate updates to the analysis by both the investigators and third
parties. This processing code will be written such that raw data is downloaded
automatically as needed. Versions  or ``releases'' of the Git repository will be
uploaded to figshare such that they will have their own DOI, and can be cited in
publications appropriately. A report detailing the experiment and the results
will be written and included as part of the repository.

\renewcommand{\bibname}{References}
\bibliography{../../../../../Library/Library}
\bibliographystyle{plain}

\end{document}